\begin{trapbox}
	\begin{description}[style=nextline, leftmargin=2.8em, labelsep=0.5em, itemsep=0.35em]
		\item[\textbf{\textcolor{CTrap}{易错点一（漏算绝对值）：}}]
		      误认为 $|\lambda\vec{a}| = \lambda|\vec{a}|$，直接略去绝对值。
		\item[\textbf{\textcolor{CTrap}{正确理解（准确模长公式）：}}]
		      $|\lambda\vec{a}| = |\lambda|\,|\vec{a}|$，其中 $\lambda$ 必须取绝对值。
		\item[\textbf{\textcolor{CTrap}{错因分析（忽略符号）：}}]
		      当 $\lambda<0$ 时，$\lambda|\vec{a}|$ 为负，而模长必须非负，故需写成 $|\lambda|\,|\vec{a}|$。
	\end{description}

	\medskip
	\begin{description}[style=nextline, leftmargin=2.8em, labelsep=0.5em, itemsep=0.35em]
		\item[\textbf{\textcolor{CTrap}{易错点二（零向量特例）：}}]
		      忽略 $\lambda=0$ 的特殊情况，认为 $\lambda\vec{a}$ 与 $\vec{a}$ 一定同向或反向。
		\item[\textbf{\textcolor{CTrap}{正确理解（零向量处理）：}}]
		      当 $\lambda=0$ 时，$\lambda\vec{a}=\vec{0}$，零向量没有确定方向，不能说与 $\vec{a}$ 同向或反向。
		\item[\textbf{\textcolor{CTrap}{错因分析（方向无定义）：}}]
		      方向关系只适合讨论非零向量；零向量只能单独处理。
	\end{description}

	\medskip
	\begin{description}[style=nextline, leftmargin=2.8em, labelsep=0.5em, itemsep=0.35em]
		\item[\textbf{\textcolor{CTrap}{易错点三（非零前提）：}}]
		      在 $\vec{a}=\vec{0}$ 时仍然套用“方向相同、方向相反”或“唯一表示”等结论。
		\item[\textbf{\textcolor{CTrap}{正确理解（强调条件）：}}]
		      本结论讨论方向时默认 $\vec{a}\neq\vec{0}$；共线定理中的唯一表示也必须以 $\vec{a}\neq\vec{0}$ 为前提。
		\item[\textbf{\textcolor{CTrap}{错因分析（零向量歧义）：}}]
		      若 $\vec{a}=\vec{0}$，则 $\lambda\vec{a}=\vec{0}$，此时方向无定义，且不能用 $\vec{b}=\lambda\vec{a}$ 唯一表示任意共线向量。
	\end{description}
\end{trapbox}